\section{Searching for an augmenting path}\label{ch08:sec:search}
Let $M$ be a matching and let $x$ be an unmatched student. How can we find an augmenting path that starts at $x$, or find out that there is none? Such a path leaves $x$ along an unmatched edge to some task. If that task is unused, the path can stop there. If it is used, the path must go on, and since the steps alternate, its next step uses a matched edge; the only one at a used task leads to its owner. From there the path continues along an unmatched edge, and so on. So we explore the graph with two kinds of moves:
\begin{itemize}
\item from a student $s$, move to any task $t$ for which $s$--$t$ is an unmatched edge;
\item from a used task $t$, move to its owner.
\end{itemize}

\subsection*{The search}
The search keeps track of which vertices it has \emph{reached}, and records for each newly reached vertex the vertex it came from, its \emph{predecessor}. To \emph{inspect} a reached student $s$, go through the unmatched edges $s$--$t$ at $s$ one at a time, and for each of them do the following.
\begin{itemize}
\item If the task $t$ has already been reached, go on to the next edge.
\item Otherwise, mark $t$ as reached, with predecessor $s$. If $t$ is unused, stop: the search has \emph{succeeded}. If $t$ is used, also mark its owner as reached, with predecessor $t$, and go on to the next edge.
\end{itemize}
The search itself runs as follows.
\begin{enumerate}
\item Mark $x$ as reached. We call $x$ the \emph{root} of the search; it has no predecessor.
\item Choose a reached student that has not been inspected yet, and inspect it.
\item Repeat step~2 until the search succeeds, or until every reached student has been inspected. In the second case the search has \emph{failed}.
\end{enumerate}
No vertex is marked twice. For a task this is built into the rule. The root $x$ owns no task, since it is unmatched, so it is never marked again. Any other student is marked only as the owner of a task being reached for the first time; since a student owns only one task, an earlier marking of that student would have been as the owner of the same task, which would then already have been reached. So each student is inspected at most once, each inspection looks at finitely many edges, and there are finitely many students, so the search always ends.

Suppose the search succeeds at an unused task $t$. Follow the predecessors backward from $t$: the predecessor of a task is a student, and the predecessor of a student other than $x$ is the task it holds. Every vertex was reached later than its predecessor, so no vertex is met twice; as there are finitely many vertices, the backward walk stops, and it can stop only at $x$, the one reached vertex without a predecessor. Read forward, it is a path from $x$ to $t$. Every step into a task used an unmatched edge and every step into a student a matched edge, so the path is alternating. Its ends are unmatched, and it has at least one step, because $x$ is a student and $t$ a task. So it is an augmenting path, and Lemma~\ref{thm:augment} lets us flip along it.

The rules do not say which uninspected student to choose next. Inspecting students in the order in which they were reached is called \emph{breadth-first search}, but nothing in this chapter depends on the order.

\subsection*{The whole procedure}
Repeating the search gives a procedure for the assignment problem:
\par\noindent\begin{minipage}{\linewidth}
\begin{lstlisting}
Start with the empty matching M.
While some student is unmatched:
    choose an unmatched student x and search from x;
    if the search succeeds:
        flip M along the augmenting path it found;
    if the search fails:
        stop, and report the set A of students it reached.
If the loop ends because every student is matched, report M.
\end{lstlisting}
\end{minipage}\par
As shown at the end of Section~\ref{ch08:sec:augment}, the procedure makes at most $|L|$ flips. When a search fails, the procedure reports a set of students, and Section~\ref{ch08:sec:shortage} shows that this set proves that no matching saturates $L$.

\subsection*{What a failed search guarantees}
When the search fails, the reached vertices have two properties.
\begin{enumerate}
\item[(C1)] If $s$ is a reached student and $s$--$t$ is an unmatched edge, then $t$ has been reached.
\item[(C2)] Every reached task is used, and its owner has been reached.
\end{enumerate}
Property~(C1) holds because a failed search has inspected every reached student, and inspecting $s$ follows every unmatched edge at $s$. Property~(C2) holds because reaching an unused task would have ended the search with success, and because the owner of a used task is marked as reached at the moment the task is.

We call a search \emph{complete} if it stops only when it has found an unused task or has inspected every reached student; the search of this section is complete. A search that gives up earlier, say when its first route is blocked, can stop while (C1) is still false, and then its failure tells us nothing: an augmenting path might use an edge it never tried.

When a complete search fails, no augmenting path starts at~$x$. Suppose there were one, $x=v_0,v_1,\ldots,v_k$. Its vertices alternate between students and tasks, starting with $x$, and its edges alternate, starting with an unmatched one (Section~\ref{ch08:sec:augment}), so the edge $v_iv_{i+1}$ is unmatched when $i$ is even, that is, when $v_i$ is a student, and matched when $v_i$ is a task. By induction on $i$, every $v_i$ has been reached. Certainly $v_0=x$ has. If a reached $v_i$ is a student, the edge $v_iv_{i+1}$ is unmatched, so $v_{i+1}$ has been reached by~(C1). If a reached $v_i$ is a task, then by~(C2) it is used and its owner has been reached; the matched edge $v_iv_{i+1}$ can only join it to its owner, so $v_{i+1}$ is that owner. Hence the last vertex $v_k$, a task, has been reached, and by~(C2) it is used, whereas the last vertex of an augmenting path is unused. This contradiction proves the claim.
